% 第20章 研究现状、遗留课题与未来路线图
\section{研究现状、遗留课题与未来路线图}\label{sec:future}

\subsection{当前研究现状}

把六个方法族横过来看，能看到一条清晰的分层：\textbf{监督学习已是常规工具，深度序列模型到达平台期，强化学习在摩擦相关环节站稳、在选股环节证据薄弱，大模型仍在爆炸期但证据等级最低，生成式模型从“造数据”走向“造环境”，图神经网络从狂热回归理性。}与之形成反差的是评测层：方法在进步，可比性却在退步——因为每出一族新方法，就多一批自定口径、缺成本、缺审计的结果，跨论文的可比性反而被稀释。

如果用一句话概括现状：\textbf{我们缺的不是新模型，是能把新模型和旧模型放在一起公平比较的公共基础设施。}

\subsection{遗留课题}

\begin{enumerate}
  \item \textbf{泄漏难排除下的 LLM 评估。}LLM 的知识截止日使“样本外”这一概念本身失效，而前向样本需要时间积累。在积累到足够前向数据之前，如何设计可信的代理评测（如强制使用截止日后数据、语料排除证明、或合成对照），是当前最关键的方法论缺口。
  \item \textbf{摩擦环节的公开基准缺失。}执行、做市、另类数据几乎无可复现基准，导致越关键、越靠近净收益的环节，证据越薄。
  \item \textbf{试验次数不可完全观测下的多重检验。}真实的超参数搜索与特征尝试次数，研究者自己也未必全知，DSR/PBO 需要的 $N$ 因此难以精确。如何在这种不确定下给出稳健的显著性判断，仍开放。
  \item \textbf{非平稳下的持续学习。}现有模型多假设“学到即固定”，而市场制度会漂移。如何在有限前向反馈下做在线适应而不被噪声带偏，是把实验室结果推向长期实盘的必经问题。
  \item \textbf{端到端与可诊断性的张力。}端到端能协调环节冲突、提升上限，却牺牲可诊断性，而交易恰恰是出错代价极高的场景。
  \item \textbf{容量与拥挤度的联合估计。}策略能赚多少、能装多少资金、别人用了它还剩多少，三者纠缠在一起，现有研究很少同时处理。
  \item \textbf{合成数据能否用于风控。}第 \ref{sec:gen} 章指出生成模型无法外推未发生事件，那么用于压力测试的合成数据，其结论可信边界在哪，需要专门研究。
\end{enumerate}

\subsection{研究路线图（Phase 0--4）}

路线图按“先立尺子、再比方法、再验前沿、最后沉淀”的顺序设计，每阶段带交付物与验证门（G1--G9）。与第 \ref{sec:repro_protocol} 章 P0--P3 的关系是：P 级是\textbf{方法纪律}，Phase 是\textbf{工程路线}，二者一一对应但阶段更长。

\begin{table}[htbp]
\centering
\caption{研究路线图：阶段、交付物与验证门}
\label{tab:roadmap}
\begin{tabular}{@{}>{\raggedright\arraybackslash}p{1.5cm}>{\raggedright\arraybackslash}p{4.6cm}>{\raggedright\arraybackslash}p{4.6cm}>{\raggedright\arraybackslash}p{3.6cm}@{}}
\toprule
阶段 & 目标 & 交付物 & 验证门 \\
\midrule
Phase 0 基建 & 立尺子 & 时点数据快照（含校验和）、三档成本层、评测容器（含 DSR/PBO）、试验账本模板 & G1 时点审计通过；G2 成本层数值可复现 \\
\midrule
Phase 1 基准复现 & 对齐已知结论 & 复现 Qlib 选股基准、经典动量/价值因子、VWAP 执行基线 & G3 复现数值落在容差内；G4 确定“最强简单基线” \\
\midrule
Phase 2 前沿对照 & 公平比较 & 在统一口径下对照监督/深度/RL 方法与基线，输出增量归因表 & G5 增量归因完成（模型 vs 设定）；G6 DSR 与 PBO 双通过 \\
\midrule
Phase 3 前沿验证 & 检验高风险方向 & LLM 前向样本评测、生成式压力测试、跨市场泛化实验 & G7 泄漏排除证明成立；G8 跨 $\geq 3$ 市场 OOS 同号 \\
\midrule
Phase 4 沉淀 & 公共资产 & 开源协议与基准、可复现排行榜、负面结果库 & G9 第三方在无作者协助下复现同一结论 \\
\bottomrule
\end{tabular}
\end{table}

\textbf{为什么把 Phase 0 放在最前}：因为如果数据、成本与评测三样不统一，后面所有比较都是沙上建塔。这也是本文与前人综述最大的不同——它不推荐你去追最新的模型，而推荐你先把地基打平。

\textbf{为什么单列负面结果库}：金融机器学习领域最稀缺的不是成功案例，而是“这个方向我试过、在统一口径下不成立”的记录。负面结果的公开，能直接降低整个领域重复踩坑的成本，它的公共价值不低于任何一篇正结果论文。

本章要记住的一句话是：\textbf{这个领域下一步最大的进展，可能不是某个模型刷出了更高的夏普，而是一套让所有人的夏普可信地接受比较的公共协议——先把尺子造出来，比抢着去量更有意义。}
